In this section, we provide detailed implementation information for our quantitative evaluation experiments. 
In section~\ref{ovssimpl}, we present the evaluation details for open-vocabulary semantic segmentation (OVSS). In section~\ref{linearimpl}, we describe the evaluation details for linear probe based semantic segmentation and monocular depth estimation.

\subsection{Implementation Details of Open-vocabulary Semantic Segmentation.}
\label{ovssimpl}
We follow SCLIP and NACLIP in the setup for the open-vocabulary semantic segmentation evaluation. For fairness, we utilize the same parameters for all models. We resize input images such that the shorter side is scaled to a specific resolution, while maintaining the original aspect ratio for the longer side. Additionally, we set fixed crop sizes and strides during evaluation. All evaluation parameters are summarized in Table~\ref{tab:resize}, while all other settings follow their default configurations.

% Please add the following required packages to your document preamble:
% \usepackage{graphicx}
\renewcommand{\arraystretch}{1.2}
\begin{table*}[t]
\centering
\caption{\textbf{Dataset Specific Details for Open-vocabulary Semantic Segmentation.} We list the per-dataset resolution, crop size, and stride used for each dataset. We maintain the same settings for all methods within a given dataset.}
\vspace{.1cm}
\label{tab:resize}
\resizebox{\columnwidth}{!}{%
\begin{tabular}{l|cccccccc}
\hline
Dataset   & VOC21       & PC 60       & Object      & VOC20       & PC 59       & Stuff       & City        & ADE         \\ \hline
Resize resolution    & 448 & 448 & 336 & 336 & 448 & 448 & 560 & 448 \\
Crop size & 336         & 336         & 336         & 336         & 336         & 336         & 224         & 336         \\
Stride    & 112         & 112         & 112         & 112         & 112         & 112         & 112         & 112         \\ \hline
\end{tabular}%
}
\end{table*}

\subsection{Implementation Details of Linear Probe Based Evalution.}
\label{linearimpl}
Our linear probe evaluation follows prior works (Vision Transformers Need Registers, Denoising Vision Transformers), where a linear layer is trained as a decoding head to predict pixel-wise segmentation or depth logits.

\textbf{Semantic Segmentation.} We extract the final output features from the frozen backbone and, if applicable, pass them through the denoiser (for the DVT baseline). A single learnable linear layer is then trained to predict the segmentation logits. For CLIP-based models, both training and testing images are resized to (448, 448), while for DINOv2, the images are resized to (518, 518).

\textbf{Monocular Depth Estimation.}
Similar to semantic segmentation, we extract features from the backbone, and pass them through the denoiser if applicable. Following the method in DVT and DINOv2, we then append the \textbf{[CLS]} token to each patch token to enrich the feature representations for all methods. 
A linear layer is trained using SigLoss and gradient loss (scaled by a factor of 0.5) to predict depth values into 256 uniformly distributed bins. We adopt DVT's learning rate of 5e-3 for all experiments. 
